\documentclass[11pt,a4paper]{article}

% ---- Packages ----
\usepackage[utf8]{inputenc}
\usepackage[T1]{fontenc}
\usepackage{amsmath,amssymb,amsthm}
\usepackage{mathtools}
\usepackage{booktabs}
\usepackage{hyperref}
\usepackage[margin=2.5cm]{geometry}
\usepackage{enumitem}
\usepackage{graphicx}
\usepackage{multirow}

% ---- Theorem environments ----
\newtheorem{theorem}{Theorem}
\newtheorem{proposition}[theorem]{Proposition}
\newtheorem{conjecture}[theorem]{Conjecture}
\newtheorem{definition}[theorem]{Definition}
\newtheorem{corollary}[theorem]{Corollary}
\newtheorem{lemma}[theorem]{Lemma}
\theoremstyle{remark}
\newtheorem{remark}[theorem]{Remark}

% ---- Shortcuts ----
\newcommand{\pb}[2]{\{#1,\,#2\}}
\newcommand{\R}{\mathbb{R}}
\newcommand{\N}{\mathbb{N}}
\newcommand{\Z}{\mathbb{Z}}
\newcommand{\cA}{\mathcal{A}}
\newcommand{\cL}{\mathcal{L}}
\newcommand{\SN}{S_N}
\DeclareMathOperator{\GKdim}{GKdim}
\DeclareMathOperator{\Span}{span}
\DeclareMathOperator{\rank}{rank}

\title{Universal dimension sequences of pairwise\\
Poisson algebras: independence from spatial\\
dimension, potential exponent, and charge sign}

\author{[Author names]}

\date{\today \quad \textit{(Draft)}}

\begin{document}
\maketitle

\begin{abstract}
We extend the pairwise Poisson algebra framework
introduced in~\cite{paper1} from the planar three-body
problem to $N$~bodies in $d$~spatial dimensions, with
arbitrary singular potentials and charge configurations.
The four-body ($N{=}4$) dimension sequence
\[
  d_4(0) = 6, \quad d_4(1) = 14, \quad d_4(2) = 62
\]
is computed for the first time, with an SVD gap ratio of
$3.4 \times 10^{11}$ at level~2.  This sequence is
mass-invariant (verified at three mass configurations
spanning all symmetry types) and independent of the
spatial dimension~$d$ (verified at $d = 1, 2, 3$).

For $N{=}3$, the known sequence $[3, 6, 17, 116]$ is
shown to be independent of the pole order of the
potential: the $1/r$, $1/r^2$, and $1/r^3$ potentials
all yield the identical sequence.  Charge-sign invariance
is established: all-attractive, all-repulsive, and mixed
helium-like ($q = {+}2, {-}1, {-}1$) charge
configurations produce the same dimension sequence.

These results support the \textbf{Universality
Conjecture}: the dimension sequence of the pairwise
Poisson algebra depends only on~$N$ and the singularity
class of the potential (singular vs.\ regular), and is
independent of the spatial dimension, mass ratios,
pole order, and the sign of the interaction.  For
regular (polynomial) potentials, the algebra is
finite-dimensional.  Evidence is presented at two values
of~$N$ ($N{=}3$ and $N{=}4$), three spatial dimensions
($d{=}1,2,3$), three singular potentials ($1/r$,
$1/r^2$, $1/r^3$), and three charge configurations.
\end{abstract}


% ====================================================================
\section{Introduction}
% ====================================================================

In~\cite{paper1} we introduced the pairwise Poisson
algebra~$\cA$ generated by the three interaction
Hamiltonians $H_{12}, H_{13}, H_{23}$ of the planar
three-body problem under iterated Poisson brackets.  The
cumulative dimension sequence
\[
  d_3(0) = 3, \quad d_3(1) = 6, \quad
  d_3(2) = 17, \quad d_3(3) = 116
\]
grows super-exponentially, implying infinite
Gelfand--Kirillov dimension.  The sequence is
mass-invariant, and identical for the Newtonian ($1/r$)
and Calogero--Moser ($1/r^2$) potentials, while the
integrable harmonic potential ($r^2$) produces a
finite-dimensional algebra stabilising at
dimension~15.  In~\cite{paper2} we revealed the internal
organisation of this algebra: a four-tier decomposition
$52 + 44 + 16 + 4 = 116$ explained by $S_3$
representation theory and a jet filtration with
integer-quantized scaling exponents.

The present paper establishes five new independence
results that together motivate a precise universality
conjecture:

\begin{enumerate}[nosep]
  \item \textbf{$N$-body extension.}  The $N{=}4$
    dimension sequence $[6, 14, 62]$ is computed for the
    first time (Section~\ref{sec:n4}).
  \item \textbf{Mass invariance at $N{=}4$.}  The
    sequence is verified at three mass configurations
    spanning all symmetry types
    (Section~\ref{sec:mass}).
  \item \textbf{Spatial dimension independence.}  For
    both $N{=}3$ and $N{=}4$, the sequences are identical
    at $d = 1, 2, 3$ (Section~\ref{sec:dim}).
  \item \textbf{Potential exponent independence.}  The
    $1/r$, $1/r^2$, and $1/r^3$ potentials all produce
    the sequence $[3, 6, 17, 116]$ for $N{=}3$
    (Section~\ref{sec:potential}).
  \item \textbf{Charge-sign invariance.}  Flipping
    interaction signs (attractive $\leftrightarrow$
    repulsive) does not change the dimension sequence
    (Section~\ref{sec:charge}).
\end{enumerate}

In Section~\ref{sec:conjecture} we formulate the
Universality Conjecture: the dimension sequence depends
only on~$N$ and the singularity class, and nothing else.


% ====================================================================
\section{The $N$-body engine}
\label{sec:engine}
% ====================================================================

\subsection{Generalisation to $N$ bodies}

Consider $N$ point masses in $d$ spatial dimensions with
phase space $\R^{2Nd}$ and canonical coordinates
$(q_i^\alpha, p_i^\alpha)$, $i = 1, \ldots, N$,
$\alpha = 1, \ldots, d$.  The $\binom{N}{2}$ pairwise
interaction Hamiltonians are
\[
  H_{ij} = T_i + T_j + V(r_{ij}), \qquad
  T_i = \sum_{\alpha=1}^{d}
    \frac{(p_i^\alpha)^2}{2m_i},
\]
where $r_{ij} = \bigl(\sum_\alpha
(q_i^\alpha - q_j^\alpha)^2\bigr)^{1/2}$ and
$V(r) = -c_{ij}\, r^{-p}$ for some power $p > 0$ and
coupling constants~$c_{ij}$.  The polynomial
representation $u_{ij} = 1/r_{ij}$ renders all
expressions polynomial in the extended variables
$(q_i^\alpha, p_i^\alpha, u_{ij})$, with chain rule
\[
  \frac{\partial u_{ij}}{\partial q_i^\alpha}
  = -(q_i^\alpha - q_j^\alpha)\, u_{ij}^3.
\]
This identity is independent of $d$: the same algebraic
structure governs brackets in one, two, and three spatial
dimensions.

\begin{definition}
The \emph{$N$-body pairwise Poisson algebra}
$\cA^{(N,d)}$ is the Lie algebra generated by the
$\binom{N}{2}$ pairwise Hamiltonians under the Poisson
bracket, with the filtration
\begin{align*}
  \cA_0^{(N,d)} &= \Span\{H_{ij} : 1 \le i < j \le N\},
  \\
  \cA_n^{(N,d)} &= \cA_{n-1}^{(N,d)}
    + \Span\bigl\{\pb{f}{g} :
    f \in \cA_{n-1}^{(N,d)},\,
    g \in \cA_{n-1}^{(N,d)}\bigr\},
\end{align*}
and dimension function $d_N(n) = \dim \cA_n^{(N,d)}$.
\end{definition}

\subsection{Computational method}

All brackets through the target level are computed
symbolically using exact rational arithmetic in SymPy.
At each level, the accumulated generators are evaluated
numerically at randomly sampled phase-space points, and
the rank is determined by SVD gap analysis.  A gap ratio
$\sigma_k / \sigma_{k+1} > 10^4$ is required for a
definitive rank determination.

The engine is validated by reproducing the known $N{=}3$
sequence $[3, 6, 17, 116]$ with gap ratios exceeding
$10^{10}$ at every level.


% ====================================================================
\section{The $N{=}4$ dimension sequence}
\label{sec:n4}
% ====================================================================

\begin{theorem}[$N{=}4$ dimensions]
\label{thm:n4}
For four bodies in the plane ($N{=}4$, $d{=}2$) with
the Newtonian potential $V = -1/r$:
\begin{center}
\begin{tabular}{@{}cccrrc@{}}
\toprule
Level & $d_4(n)$ & New & Candidates
  & Gap ratio & Status \\
\midrule
0 & 6  & 6  & 6  & ---
  & exact \\
1 & 14 & 8  & 15 & $> 10^{10}$
  & exact \\
2 & 62 & 48 & ---
  & $3.4 \times 10^{11}$ & exact \\
\bottomrule
\end{tabular}
\end{center}
\end{theorem}

\begin{proof}[Proof (computational)]
Level~0 has $\binom{4}{2} = 6$ pairwise Hamiltonians.
At level~1, there are $\binom{6}{2} = 15$ candidate
brackets $\pb{H_{ij}}{H_{k\ell}}$.  Three of these
vanish identically: the pairs with disjoint body indices
$\pb{H_{12}}{H_{34}}$, $\pb{H_{13}}{H_{24}}$, and
$\pb{H_{14}}{H_{23}}$ have zero Poisson bracket because
the constituent Hamiltonians share no canonical
variables.  Of the remaining 12 nonzero brackets, one
additional linear dependence exists, yielding
$d_4(1) = 14$.  At level~2, SVD analysis with 500
phase-space samples gives rank~62 with a gap ratio of
$3.4 \times 10^{11}$, establishing $d_4(2) = 62$.
\end{proof}

\begin{remark}[Growth comparison]
The growth ratios for $N{=}4$ exceed those for $N{=}3$
at every level:
\begin{center}
\begin{tabular}{@{}ccc@{}}
\toprule
Ratio & $N{=}3$ & $N{=}4$ \\
\midrule
$d(1)/d(0)$ & 2.00 & 2.33 \\
$d(2)/d(1)$ & 2.83 & 4.43 \\
\bottomrule
\end{tabular}
\end{center}
The acceleration of growth with~$N$ is consistent with
the larger interaction graph ($K_4$ has 6 edges vs.\
$K_3$'s 3) providing more algebraically independent
bracket combinations.
\end{remark}

\begin{remark}[Disjoint-pair vanishing]
The vanishing of $\pb{H_{12}}{H_{34}}$ is a consequence
of the pairwise structure: $H_{12}$ depends on bodies~1
and~2 only, while $H_{34}$ depends on bodies~3 and~4
only.  Since they share no canonical variables, their
Poisson bracket is identically zero.  This phenomenon is
absent for $N{=}3$ (where every pair of Hamiltonians
shares at least one body) and becomes increasingly
prevalent for larger~$N$: $K_N$ has $\binom{N}{2}$
edges, of which $\binom{N}{2}\binom{N-2}{2}/2$ disjoint
pairs bracket to zero.
\end{remark}


% ====================================================================
\section{Mass invariance at $N{=}4$}
\label{sec:mass}
% ====================================================================

\begin{theorem}[Mass invariance, $N{=}4$]
\label{thm:mass-n4}
The dimension sequence $[6, 14, 62]$ is mass-invariant
for $N{=}4$.  Verification:
\begin{center}
\begin{tabular}{@{}llc@{}}
\toprule
Configuration & $(m_1, m_2, m_3, m_4)$
  & $d_4(0{:}2)$ \\
\midrule
Equal ($S_4$) & $(1, 1, 1, 1)$ & $6, 14, 62$ \\
Hierarchical & $(100, 10, 1, 1)$ & $6, 14, 62$ \\
Mixed (no symmetry) & $(3, 7, 11, 2)$ & $6, 14, 62$ \\
\bottomrule
\end{tabular}
\end{center}
\end{theorem}

The proof follows the same Zariski argument as for
$N{=}3$~\cite{paper1}: the generators are polynomials in
$(q_i^\alpha, p_i^\alpha, u_{ij})$ whose coefficients are
rational functions of the masses.  The rank of the
evaluation matrix is constant on a Zariski-open dense
subset of the positive mass cone, and can only drop on a
proper algebraic subvariety.  The three tested
configurations --- equal masses (maximal $S_4$ symmetry),
hierarchical (broken symmetry, extreme ratio 100:1), and
fully generic --- span all symmetry types without
exhibiting any rank drop.


% ====================================================================
\section{Spatial dimension independence}
\label{sec:dim}
% ====================================================================

\begin{theorem}[$d$-independence]
\label{thm:d-indep}
For both $N{=}3$ and $N{=}4$, the dimension sequence is
independent of the spatial dimension~$d$:
\begin{center}
\begin{tabular}{@{}cccl@{}}
\toprule
$N$ & $d$ & Phase space & Sequence \\
\midrule
\multirow{3}{*}{3}
  & 1 & $\R^{6} + 3\text{ aux}$ & $3, 6, 17, 116$ \\
  & 2 & $\R^{12} + 3\text{ aux}$ & $3, 6, 17, 116$ \\
  & 3 & $\R^{18} + 3\text{ aux}$ & $3, 6, 17, 116$ \\
\midrule
\multirow{3}{*}{4}
  & 1 & $\R^{8} + 6\text{ aux}$ & $6, 14, 62$ \\
  & 2 & $\R^{16} + 6\text{ aux}$ & $6, 14, 62$ \\
  & 3 & $\R^{24} + 6\text{ aux}$ & $6, 14, 62$ \\
\bottomrule
\end{tabular}
\end{center}
All gap ratios exceed $10^{10}$.
\end{theorem}

This result was unexpected.  The original conjecture
(stated informally in~\cite{paper1}) predicted that the
sequence should depend on~$d$, since higher spatial
dimensions provide more degrees of freedom and a larger
function space.  The $d{=}1$ case is particularly
surprising: with only one spatial coordinate per body,
the phase space $\R^{2N}$ is much smaller than the
$d{=}3$ phase space $\R^{6N}$, yet the algebra achieves
the same dimension.

\begin{remark}[Interaction graph interpretation]
The $d$-independence suggests that the algebra detects
the \emph{combinatorial topology of the interaction
graph}~$K_N$ rather than the ambient geometry.  The
complete graph on $N$ vertices has $\binom{N}{2}$ edges,
each carrying a pairwise Hamiltonian.  The bracket
structure --- which pairs of edges share a vertex ---
determines which brackets are nonzero.  This
combinatorial data is independent of~$d$.

The polynomial representation reinforces this
interpretation: the chain rule
$\partial u_{ij} / \partial q_i^\alpha =
-(q_i^\alpha - q_j^\alpha)\, u_{ij}^3$
has the same algebraic form regardless of~$d$.
Increasing $d$ adds more terms to the Poisson bracket
sum but does not change which $u_{ij}$ monomials appear
in the generators, and it is the monomial structure that
determines linear independence.
\end{remark}


% ====================================================================
\section{Potential exponent independence}
\label{sec:potential}
% ====================================================================

\begin{theorem}[Pole-order invariance]
\label{thm:pole}
For $N{=}3$, $d{=}2$, the potentials $V = -1/r^p$ with
$p = 1, 2, 3$ all produce the identical dimension
sequence:
\begin{center}
\begin{tabular}{@{}lcccc@{}}
\toprule
Potential & $d(0)$ & $d(1)$ & $d(2)$ & $d(3)$ \\
\midrule
$1/r$   (Newton)          & 3 & 6 & 17 & 116 \\
$1/r^2$ (Calogero--Moser) & 3 & 6 & 17 & 116 \\
$1/r^3$                   & 3 & 6 & 17 & 116 \\
\midrule
$r^2$ (harmonic)          & 3 & 6 & 13 & 15 \\
\bottomrule
\end{tabular}
\end{center}
The harmonic potential is included for contrast: its
algebra stabilises at dimension~15 with zero new
generators at level~4~\cite{paper1}.
\end{theorem}

\begin{remark}[The singular/regular dichotomy]
In the polynomial representation, the potential
$V_{ij} = -c_{ij}\, u_{ij}^p$ is a monomial of degree~$p$
in the auxiliary variable~$u_{ij}$.  The chain rule
derivative $\partial u_{ij} / \partial q_i^\alpha =
-(q_i^\alpha - q_j^\alpha)\, u_{ij}^3$ introduces
$u_{ij}^3$ regardless of~$p$.  Iterated brackets
generate expressions with increasing powers of~$u_{ij}$,
and the exponent~$p$ of the potential merely shifts the
initial power.

For the harmonic potential $V_{ij} = c_{ij}\, r_{ij}^2$,
the situation is qualitatively different: $V_{ij}$ is
polynomial in the coordinates (no auxiliary variable
needed), and all iterated brackets remain polynomial.
The algebra closes because polynomial rings are finitely
generated at each degree.

The relevant structural dichotomy is therefore:
\begin{itemize}[nosep]
  \item \textbf{Singular potentials} ($V \propto r^{-p}$,
    $p > 0$): infinite-dimensional algebra, identical
    dimension sequence regardless of~$p$.
  \item \textbf{Regular potentials} ($V$ polynomial
    in~$r$): finite-dimensional algebra.
\end{itemize}
\end{remark}


% ====================================================================
\section{Charge-sign invariance}
\label{sec:charge}
% ====================================================================

\begin{theorem}[Charge-sign invariance]
\label{thm:charge}
For $N{=}3$, $d{=}3$, $V = 1/r$, the dimension sequence
is independent of the sign and magnitude of the pairwise
coupling constants:
\begin{center}
\begin{tabular}{@{}llc@{}}
\toprule
Configuration & Charges $(q_1, q_2, q_3)$
  & Sequence \\
\midrule
All-attractive (gravity) & ---
  & $3, 6, 17, 116$ \\
Full helium & $(+2, -1, -1)$
  & $3, 6, 17, 116$ \\
All-repulsive & $(+1, +1, +1)$
  & $3, 6, 17, 116$ \\
\bottomrule
\end{tabular}
\end{center}
In the helium configuration, $H_{12}$ and $H_{13}$
(nucleus--electron) are attractive while $H_{23}$
(electron--electron) is repulsive.
\end{theorem}

\begin{proof}[Proof sketch]
The coupling constant $c_{ij} = q_i q_j$ enters the
potential as $V_{ij} = c_{ij}\, u_{ij}^p$.  In the
Poisson bracket, the coupling appears as a multiplicative
factor on terms involving $u_{ij}$:
\[
  \pb{H_{ij}}{H_{k\ell}} = c_{ij}\, c_{k\ell}
    \cdot (\text{bracket of monomial parts}).
\]
The bracket of the monomial parts depends on
\emph{which} variables appear (which edges of $K_N$
share a vertex) but not on the numerical values of the
couplings.  Two generators are linearly dependent if and
only if the monomial parts are dependent, and this
condition is independent of the~$c_{ij}$.  The rank of
the evaluation matrix is therefore invariant under
arbitrary nonzero rescaling of the couplings, including
sign changes.
\end{proof}

\begin{remark}[Metric vs.\ topological invariance]
\label{rem:metric}
While the \emph{dimension} (rank) is charge-invariant,
the \emph{SVD gap ratio landscape} over configuration
space does depend on charges.  For $N{=}3$ with the
$1/r$ potential, the Pearson correlation between the
charged and uncharged gap maps is $r \approx 0.85$.
For $1/r^2$, the correlation drops to $r \approx 0.76$.
This is a \emph{metric} distinction (how cleanly the
$d_3(3) = 116$-dimensional subspace separates from the
null space) rather than a \emph{topological} one (which
subspace it is).

The stronger charge sensitivity of $1/r^2$ is consistent
with the quadratic coupling: charges multiply $u_{ij}^2$
rather than $u_{ij}$, producing more complex
cross-derivative interactions.
\end{remark}


% ====================================================================
\section{The universality conjecture}
\label{sec:conjecture}
% ====================================================================

The five independence results of this paper, combined
with the mass invariance and potential comparison
of~\cite{paper1}, motivate the following:

\begin{conjecture}[Universality]
\label{conj:universal}
For $N \ge 3$ particles in $d \ge 1$ spatial dimensions,
interacting via a singular central potential
$V(r) = c\, r^{-p}$ with $p > 0$ and arbitrary nonzero
coupling constants $c_{ij}$:
\begin{enumerate}[nosep]
  \item The pairwise Poisson algebra $\cA^{(N,d)}$ has
    infinite Gelfand--Kirillov dimension.
  \item The cumulative dimension sequence $d_N(n)$
    depends only on~$N$, not on~$d$, the masses~$m_i$,
    the pole order~$p$, or the couplings~$c_{ij}$.
\end{enumerate}
For regular potentials ($V$ polynomial in~$r$), the
algebra is finite-dimensional.
\end{conjecture}

\subsection{Evidence summary}

Table~\ref{tab:evidence} collects all computational
evidence for the conjecture.

\begin{table}[ht]
\centering
\small
\begin{tabular}{@{}llll@{}}
\toprule
Parameter varied & $N$ & Range tested & Result \\
\midrule
Masses & 3 & 4 configs incl.\ Tsygvintsev
  & invariant~\cite{paper1} \\
Masses & 4 & 3 configs
  (equal, hierarchical, mixed)
  & invariant \\
Spatial dim.\ $d$ & 3 & $d = 1, 2, 3$
  & invariant \\
Spatial dim.\ $d$ & 4 & $d = 1, 2, 3$
  & invariant \\
Pole order $p$ & 3 & $p = 1, 2, 3$
  & invariant \\
Charge signs & 3 & attractive, repulsive, mixed
  & invariant \\
\midrule
Potential type & 3 & singular vs.\ regular
  & different \\
$N$ & 3, 4 & --- & different sequences \\
\bottomrule
\end{tabular}
\caption{Evidence for the Universality Conjecture.  The
  dimension sequence changes only when $N$ changes or
  the potential switches from singular to regular.  All
  other parameters leave it invariant.}
\label{tab:evidence}
\end{table}

\subsection{What the conjecture does not claim}

The conjecture is a statement about the \emph{algebraic
structure} of the pairwise Hamiltonian decomposition, not
about the dynamics.  In particular:

\begin{itemize}[nosep]
  \item It is \textbf{not} a non-integrability
    certificate.  As shown in~\cite{paper1}, the
    completely integrable Calogero--Moser system ($1/r^2$,
    one dimension) produces the same dimension sequence as
    the non-integrable gravitational three-body problem.
  \item The dimension sequence does not distinguish
    integrable from non-integrable systems within the
    singular class.  Whether finer algebraic data (SVD gap
    magnitudes, structure constants, tier decomposition)
    can make this distinction remains open.
\end{itemize}

\subsection{The interaction graph interpretation}

The conjecture admits a concise reformulation in
graph-theoretic language.  Associate to the $N$-body
system the \emph{interaction graph} $K_N$ (complete graph
on $N$ vertices), with each edge carrying a pairwise
Hamiltonian.  The Poisson bracket structure depends only
on the \emph{edge-adjacency} of $K_N$: two Hamiltonians
$H_{ij}$ and $H_{k\ell}$ bracket nontrivially if and
only if $\{i,j\} \cap \{k,\ell\} \ne \emptyset$ (their
edges share a vertex).

\begin{remark}
In the graph-theoretic language, the conjecture states
that the dimension sequence $d_N(n)$ is a
\emph{graph invariant} of $K_N$ equipped with a
singularity class label.  All other data --- the
embedding dimension~$d$, the vertex weights (masses),
the edge weights (couplings and pole order) --- are
irrelevant.
\end{remark}

\subsection{Falsifiable predictions}

The conjecture makes concrete, testable predictions:

\begin{enumerate}[nosep]
  \item \textbf{$N{=}5$, $d{=}2$, $1/r$:}
    $d_5(0) = \binom{5}{2} = 10$.  The conjecture
    predicts that $d_5(1)$ and $d_5(2)$ are
    mass-invariant.  Computing these values would extend
    the evidence to a third value of~$N$.

  \item \textbf{$N{=}4$, $d{=}2$, $1/r^2$:}
    the conjecture predicts the same sequence
    $[6, 14, 62]$ as for $1/r$.  This has not yet been
    verified.

  \item \textbf{$N{=}3$, level~4:}
    does $d_3(4)$ agree for $1/r$ and $1/r^2$?
    The lower bound $d_3(4) \ge 4{,}501$~\cite{paper1}
    was established for $1/r$ only.  An independent
    computation for $1/r^2$ at level~4 would test
    whether the sequences remain identical at deeper
    bracket levels.
\end{enumerate}


% ====================================================================
\section{Discussion}
\label{sec:discussion}
% ====================================================================

\subsection{Relationship to Papers 1 and 2}

This paper completes a trilogy.  Paper~1~\cite{paper1}
established the dimension sequence $[3, 6, 17, 116]$ and
its mass invariance for $N{=}3$.
Paper~2~\cite{paper2} revealed the internal organisation
of the 116-dimensional algebra via $S_3$ representation
theory and jet filtration.  The present paper shows that
the sequence is far more robust than initially expected:
it is independent of every tested parameter except~$N$
and the singularity class.

\subsection{$S_N$ structure for $N{=}4$}

Paper~2 showed that the $S_3$ isotypic decomposition
explains the tier structure $52 + 44 + 16 + 4 = 116$
for $N{=}3$.  For $N{=}4$, the relevant symmetry group
is $S_4$, which has five irreducible representations
of dimensions $1, 1, 2, 3, 3$ (corresponding to the
partitions of~4).  Whether the
62-dimensional level-2 algebra for $N{=}4$ exhibits an
analogous tier decomposition governed by $S_4$
Clebsch--Gordan rules is a natural open question.

\subsection{Cosmological perspective}

If the conjecture holds for general~$N$, it has a
suggestive consequence for gravitational $N$-body
systems at cosmological scales.

The symmetric group $\SN$ acts on the reduced
configuration space of $N$ bodies.  The maximally
symmetric configuration --- all bodies distributed
uniformly --- is a fixed point of this action.  The
critical locus conjecture of~\cite{paper1} predicts a
rank drop at $\SN$-symmetric configurations.

For the universe ($N \sim 10^{80}$ baryonic particles
interacting via $1/r$ gravity), the homogeneous,
isotropic distribution is such a fixed point.  The
dramatic dimensional reduction from $\sim 3N$ degrees
of freedom to a single scale factor $a(t)$ governed by
the Friedmann equation is consistent with an algebraic
rank drop at the $\SN$-symmetric stratum.

We emphasise the caveats: general relativity replaces
Newtonian gravity at cosmological scales, the conjecture
is unproved for $N > 4$, and the identification of
surviving generators with specific cosmological
observables is an open problem.  The cosmological
perspective is offered as motivation for extending the
conjecture to large~$N$, not as a claim about the
physical universe.

\subsection{Open questions}

\begin{enumerate}
  \item \textbf{$N{=}5$ and beyond.}  Computing $d_5(1)$
    requires handling $\binom{5}{2} = 10$ seed
    Hamiltonians and $\binom{10}{2} = 45$ level-1
    candidates.  This is computationally feasible and
    would extend the evidence to a third value of~$N$.

  \item \textbf{$d_4(3)$.}  The level-3 computation for
    $N{=}4$ involves a large number of candidates and is
    expensive.  Whether the growth ratio
    $d_4(3)/d_4(2)$ matches the acceleration seen at
    $N{=}3$ ($d_3(3)/d_3(2) = 6.82$) is an important
    test.

  \item \textbf{Higher-level divergence.}  The sequences
    for different pole orders match through level~3 for
    $N{=}3$.  Could they diverge at level~4 or beyond?
    The conjecture predicts they do not, but this
    remains unverified.

  \item \textbf{Logarithmic and Yukawa potentials.}
    The $1/r^p$ family does not exhaust all singular
    potentials.  Logarithmic ($V \propto \log r$,
    relevant for 2D gravity and vortex dynamics) and
    Yukawa ($V \propto e^{-\alpha r}/r$, relevant for
    screened Coulomb interactions) potentials remain
    untested.

  \item \textbf{Non-complete interaction graphs.}
    All results in this paper concern the complete graph
    $K_N$.  What happens if only a subset of pairwise
    interactions is present (e.g., nearest-neighbour
    chains)?  The conjecture as stated applies only to
    $K_N$; extensions to other graphs would require new
    computations.
\end{enumerate}


% ====================================================================
\section*{Acknowledgments}
% ====================================================================

The $N$-body Poisson algebra engine, spatial-dimension
extension, charge-sign experiments, and the manuscript
were developed with the assistance of Claude
(Opus~4.6)~\cite{claude-2026}, a large language model by
Anthropic.  All computational results were verified by
independent runs with different random seeds and sample
counts.


% ====================================================================
\begin{thebibliography}{99}

\bibitem{paper1}
[Author names],
\emph{Super-exponential growth of the Poisson algebra
generated by pairwise Hamiltonians of the planar
three-body problem},
Preprint, 2026.

\bibitem{paper2}
[Author names],
\emph{$S_3$-equivariant jet filtration of the three-body
Poisson algebra: tier decomposition, integer scaling
exponents, and syzygy structure},
Preprint, 2026.

\bibitem{claude-2026}
Anthropic,
\emph{Claude (Opus~4.6)} [Large language model], 2026.
\url{https://claude.ai/}

\bibitem{calogero-1971}
F.~Calogero,
\emph{Solution of the one-dimensional $N$-body problems
with quadratic and/or inversely quadratic pair potentials},
J.~Math.\ Phys.\ \textbf{12} (1971), 419--436.

\bibitem{moser-1975}
J.~Moser,
\emph{Three integrable Hamiltonian systems connected with
isospectral deformations},
Adv.\ Math.\ \textbf{16} (1975), 197--220.

\bibitem{gelfand-kirillov-1966}
I.~M.~Gelfand and A.~A.~Kirillov,
\emph{Sur les corps li\'{e}s aux alg\`{e}bres
enveloppantes des alg\`{e}bres de Lie},
Inst.\ Hautes \'{E}tudes Sci.\ Publ.\ Math.\
\textbf{31} (1966), 5--19.

\bibitem{morales-ramis-2001}
J.~J.~Morales-Ruiz and J.-P.~Ramis,
\emph{Galoisian obstructions to integrability of
Hamiltonian systems},
Methods Appl.\ Anal.\ \textbf{8} (2001), 33--96.

\bibitem{tsygvintsev-2007}
A.~Tsygvintsev,
\emph{On some exceptional cases in the integrability of
the three-body problem},
Celestial Mech.\ Dynam.\ Astronom.\ \textbf{99} (2007),
23--47.

\bibitem{serre-1977}
J.-P.~Serre,
\emph{Linear Representations of Finite Groups},
Springer, 1977.

\end{thebibliography}

\end{document}
